\section{The height bound}
\label{sec:height}

The balance condition of \Cref{def:avl} constrains one node at a time, and we
now show that it constrains the whole tree. The argument runs backwards from the
question one might expect. Rather than asking how tall a tree of $n$ nodes can
become, we ask how few nodes a tree of height $h$ can contain. If even the
sparsest legal tree of height $h$ needs many nodes, then a tree with few nodes
cannot be tall, and that is the bound we want.

\begin{definition}[Minimum size]
\label{def:nh}
Let $\Nmin(h)$ denote the smallest number of nodes in any AVL tree of height $h$.
\end{definition}

\subsection{A recurrence for the minimum size}

\begin{lemma}[Minimum-size recurrence]
\label{lem:recurrence}
The quantity $\Nmin(h)$ of \Cref{def:nh} satisfies $\Nmin(0) = 1$,
$\Nmin(1) = 2$, and
\begin{equation}
  \label{eq:nrec}
  \Nmin(h) = \Nmin(h-1) + \Nmin(h-2) + 1 \qquad \text{for } h \ge 2 .
\end{equation}
\end{lemma}

\begin{proof}
A tree of height $0$ is a single node, so $\Nmin(0) = 1$. A tree of height $1$
has a root and at least one child, and one child is enough, because a balance
factor of $\pm 1$ is legal. Hence $\Nmin(1) = 2$.

Let $h \ge 2$ and let $T$ be an AVL tree of height $h$ with as few nodes as
possible. Its root has some subtree of height $h-1$, since otherwise $T$ would be
shorter. Name that subtree $A$ and the other one $B$. By \Cref{def:avl} the
height of $B$ is at least $h-2$, and making it taller than $h-2$ would only add
nodes, so a minimum-size tree takes $B$ of height exactly $h-2$. Both $A$ and $B$
are themselves AVL trees, and each must be as small as its own height permits,
because replacing either by a smaller tree of the same height would shrink $T$
without changing its height. Counting the root together with the two subtrees
gives \Cref{eq:nrec}.
\end{proof}

\Cref{fig:mintrees} draws the first five of these trees. Each one consists of the
two before it placed under a new root, which is the recurrence made visible.

\begin{figure}[t]
  \centering
  % The recursion behind the height bound, drawn as the objects it describes.
% T_h is built from T_{h-1} and T_{h-2} under one new root, so the node counts
% run 1, 2, 4, 7, 12 --- each one the sum of the previous two, plus the root.
%
% The macro doing the drawing is the same one the talk uses. It places each
% subtree by the golden ratio through Binet's formula, so the picture that
% argues for the golden ratio is itself laid out by it.
\begin{tikzpicture}[x = 1mm, y = 1mm]

  \setavlscale{3.1}{6.2}

  \foreach \h [count = \i from 0] in {0, 1, 2, 3, 4} {
    \begin{scope}[xshift = \i * 27mm]
      \resetavl
      \drawfib{\h}{0}{0}
      \node[font = \small] at (0, -33) {$T_{\h}$};
    \end{scope}
  }

  % node counts under each tree
  \foreach \n [count = \i from 0] in {1, 2, 4, 7, 12}
    \node[tlabel] at (\i * 27, -40) {$N = \n$};

  \node[tannot] at (54, -50)
    {each tree is the two before it, joined under one new root};

\end{tikzpicture}

  \caption{Minimum-size AVL trees of heights $0$ through $4$. Each tree is built
  from the two preceding trees under a fresh root, so the node counts $1, 2, 4,
  7, 12$ obey \Cref{eq:nrec}.}
  \label{fig:mintrees}
\end{figure}

The recurrence in \Cref{eq:nrec} is the Fibonacci recurrence with an extra
constant, and that constant can be absorbed.

\begin{theorem}[Closed form]
\label{thm:closedform}
For every $h \ge 0$,
\begin{equation}
  \label{eq:closedform}
  \Nmin(h) = F(h+3) - 1 ,
\end{equation}
where $F$ is the Fibonacci sequence with $F(0) = 0$ and $F(1) = 1$.
\end{theorem}

\begin{proof}
Put $M(h) = \Nmin(h) + 1$. Adding $1$ to both sides of \Cref{eq:nrec} gives
\begin{equation}
  \label{eq:mrec}
  M(h) = M(h-1) + M(h-2) ,
\end{equation}
because the two constants on the right of \Cref{eq:nrec} combine with the
original one. So the shifted sequence obeys the Fibonacci recurrence exactly,
with no remainder. Its initial values are $M(0) = 2 = F(3)$ and $M(1) = 3 =
F(4)$. Two sequences that satisfy the same second-order recurrence and agree at
two consecutive indices agree everywhere, so $M(h) = F(h+3)$ for all $h \ge 0$.
Subtracting the $1$ gives \Cref{eq:closedform}.
\end{proof}

\Cref{tab:nh} lists the resulting counts against the size of a perfect tree of
the same height. The sparsest legal tree of height $15$ holds $2583$ nodes where
a perfect tree holds $65\,535$, so the balance condition tolerates a tree filled
to under four percent of its capacity. What it does not tolerate is a tree that
is tall and nearly empty, and the table shows why: the minimum size grows by a
factor close to $1.618$ at every step.

\begin{table}[t]
  \centering
  \caption{The fewest nodes an AVL tree of height $h$ can hold, against the
  $2^{h+1}-1$ nodes of a perfect tree of that height. The ratio falls steadily,
  but $\Nmin(h)$ still grows exponentially.}
  \label{tab:nh}
  \begin{tabular}{@{}S[table-format=2.0]S[table-format=4.0]S[table-format=5.0]S[table-format=3.1]@{\qquad\qquad}S[table-format=2.0]S[table-format=4.0]S[table-format=5.0]S[table-format=2.1]@{}}
    \toprule
    {$h$} & {$\Nmin(h)$} & {perfect} & {\%} &
    {$h$} & {$\Nmin(h)$} & {perfect} & {\%} \\
    \midrule
    0 &    1 &     1 & 100.0 &  8 &   88 &   511 & 17.2 \\
    1 &    2 &     3 &  66.7 &  9 &  143 &  1023 & 14.0 \\
    2 &    4 &     7 &  57.1 & 10 &  232 &  2047 & 11.3 \\
    3 &    7 &    15 &  46.7 & 11 &  376 &  4095 &  9.2 \\
    4 &   12 &    31 &  38.7 & 12 &  609 &  8191 &  7.4 \\
    5 &   20 &    63 &  31.7 & 13 &  986 & 16383 &  6.0 \\
    6 &   33 &   127 &  26.0 & 14 & 1596 & 32767 &  4.9 \\
    7 &   54 &   255 &  21.2 & 15 & 2583 & 65535 &  3.9 \\
    \bottomrule
  \end{tabular}
\end{table}

\subsection{From the recurrence to a logarithm}

The factor $1.618$ is not a coincidence, and identifying it turns
\Cref{eq:closedform} into a statement about height.

\begin{lemma}[Growth rate]
\label{lem:growth}
Let $\varphi = (1+\sqrt{5})/2$. Then $M(h) = \Nmin(h) + 1$ satisfies $M(h) =
\Theta(\varphi^{h})$.
\end{lemma}

\begin{proof}
Look for solutions of \Cref{eq:mrec} of the form $M(h) = x^{h}$. Substituting and
dividing by $x^{h-2}$ turns the recurrence into the characteristic equation
\begin{equation}
  \label{eq:char}
  x^{2} = x + 1 ,
\end{equation}
whose roots are $\varphi = (1+\sqrt{5})/2 \approx 1.618$ and $\psi =
(1-\sqrt{5})/2 \approx -0.618$. Every solution of \Cref{eq:mrec} is a
combination $M(h) = \alpha\varphi^{h} + \beta\psi^{h}$. Since $|\psi| < 1$, the
second term tends to $0$, and the first term dominates. The initial values
$M(0) = 2$ and $M(1) = 3$ give a positive $\alpha$, so $M(h) =
\Theta(\varphi^{h})$.
\end{proof}

\Cref{lem:growth} says that the sparsest legal tree multiplies its node count by
about $1.618$ for each extra level. Node count grows exponentially in height, so
height grows logarithmically in node count, and the base of that exponential
fixes the constant.

\begin{theorem}[Height bound]
\label{thm:heightbound}
Every AVL tree with $n \ge 1$ nodes has height
\begin{equation}
  \label{eq:heightbound}
  h \;<\; \frac{1}{\lgg\varphi}\,\lgg(n+2) \;+\; \frac{\lgg 5}{2\lgg\varphi} - 3
  \;\approx\; 1.4404\,\lgg(n+2) - 1.3277 .
\end{equation}
\end{theorem}

\begin{proof}
Let $T$ be an AVL tree of height $h$ with $n$ nodes. By \Cref{def:nh} and
\Cref{thm:closedform},
\begin{equation}
  \label{eq:nge}
  n \;\ge\; \Nmin(h) \;=\; F(h+3) - 1 .
\end{equation}
Binet's formula states that $F(k) = (\varphi^{k} - \psi^{k})/\sqrt{5}$ with
$\psi = (1-\sqrt{5})/2$. Since $|\psi| < 1$, we have $|\psi^{k}| \le 1$ for every
$k \ge 0$, and therefore
\begin{equation}
  \label{eq:fiblower}
  F(k) \;\ge\; \frac{\varphi^{k} - 1}{\sqrt{5}} \;>\; \frac{\varphi^{k}}{\sqrt{5}} - 1 .
\end{equation}
Combining \Cref{eq:nge} with \Cref{eq:fiblower} at $k = h+3$ gives
$n + 2 > \varphi^{h+3}/\sqrt{5}$. Taking logarithms to base two,
\begin{equation}
  \label{eq:logstep}
  \lgg(n+2) \;>\; (h+3)\lgg\varphi - \tfrac{1}{2}\lgg 5 ,
\end{equation}
and solving \Cref{eq:logstep} for $h$ yields \Cref{eq:heightbound}.
\end{proof}

Two checks show that \Cref{eq:heightbound} is close to exact rather than merely
correct. The minimum-size tree of height $7$ holds $54$ nodes, and the bound
evaluated at $n = 54$ gives $7.04$. At $n = 10^{6}$ the bound gives $27.38$,
and the largest height any AVL tree of one million nodes can reach, read off
\Cref{eq:nrec} directly, is $27$. The bound loses less than one level across five
orders of magnitude.

Against the perfect tree the cost of the balance condition is therefore modest. A
perfect tree of $10^{6}$ nodes has height $19$, so the worst permissible AVL tree
of that size is about $44$ percent taller, and every operation on it still costs
a logarithmic number of comparisons.

\begin{corollary}[Operation cost]
\label{cor:opcost}
Searching an AVL tree of $n$ nodes costs $O(\log n)$ comparisons.
\end{corollary}

\begin{proof}
Combine \Cref{prop:searchcost} with \Cref{thm:heightbound}.
\end{proof}

\begin{remark}[Two height conventions]
\label{rem:conventions}
Sources disagree about what the height of a single node is, and the disagreement
changes every constant above. We count edges, so a single node has height $0$ and
\Cref{thm:closedform} reads $\Nmin(h) = F(h+3) - 1$. Sources that count levels
instead give a single node the height $1$ and obtain $\Nmin(h) = F(h+2) - 1$,
which shifts the additive constant in \Cref{eq:heightbound} from $-1.3277$ to
$-0.3277$. Both statements are correct under their own convention, and neither
transfers to the other. A reader comparing \Cref{eq:heightbound} against
Knuth~\cite{knuth3} or against a textbook should check which convention is in
force before concluding that one of them is wrong.
\end{remark}
